\section*{How to use these notes}
\addcontentsline{toc}{section}{How to use these notes}

These notes build the thesis model one layer at a time, starting from Bishop's textbook mixture of experts and ending with the full regime-gated model. They cover the \emph{formulation} only: what the model is, which distribution it defines, what its likelihood is, and what it predicts. How we train it in the code (Adam, early stopping, walk-forward folds) comes later.

The order is deliberate.
\begin{enumerate}
  \item \textbf{Section 1} translates Bishop's notation into ours.
  \item \textbf{Section 2 is the warm-up}: a reading guide to Chamroukhi \& Nguyen (2018) and Bishop, followed by two exercises. Exercise A1 is a mixture of experts with \emph{linear-regression} experts and a softmax gate. Exercise A2 replaces the softmax gate with a hidden Markov chain. Because the experts are linear, EM has closed-form steps, so you can carry out every derivation by hand.
  \item \textbf{Sections 3 to 5} replace the linear experts with MLPs. They explain where the softmax gate comes from, why gradient training of the likelihood behaves like EM, how backpropagation works, and how the softmax gate becomes an HMM gate.
  \item \textbf{Sections 6 to 9} assemble our model: the full generative model, the likelihood, prediction and posteriors, and a checklist of assumptions and open design choices.
  \item \textbf{Appendices A and B} give the full solutions to the two exercises. \emph{Attempt the exercises first.}
  \item \textbf{Appendix D} gives the theory of the per-date and per-row likelihoods behind Q24.
\end{enumerate}

\subsection*{Notation used throughout}

\begin{center}
\renewcommand{\arraystretch}{1.25}
\begin{tabularx}{\textwidth}{@{}lX@{}}
\toprule
Symbol & Meaning \\
\midrule
$t = 1,\dots,T$ & trading day \\
$i \in \mathcal S_t$, $N_t = |\mathcal S_t|$ & stocks in the universe on day $t$ \\
$x_{i,t}\in\R^d$ & characteristics of stock $i$, known at the close of day $t$ \\
$r_{i,t+1}$ & the \textbf{target}: the \emph{market-neutral} return of stock $i$ over $(t,t+1]$, i.e.\ its return minus that day's equal-weighted cross-sectional mean (decided 2026-10-01, Q25) \\
$m_t$ & the gate's input series: the daily log return of the CRSP value-weighted S\&P 500 index (Q23) \\
$\F_t=\sigma(m_1,\dots,m_t)$ & the gate's information at the close of day $t$: the $\sigma$-algebra generated by the market series up to $t$. The family $(\F_t)_{t\ge1}$ is a filtration (Section 1.4) \\
$\mathcal G_t\supseteq\F_t$ & all information at the close of $t$: market series, characteristics and returns up to $t$ \\
$z_t\in\{1,\dots,K\}$ & market regime on day $t$ (one-hot vector $z_{tk}$) \\
$A_{jk}$, $\rho_k$ & transition probability $p(z_{t+1}{=}k\mid z_t{=}j)$ and initial probability $p(z_1{=}k)$ \\
$\eta_t(k)=\N(m_t;\nu_k,\sigma_k^2)$ & gate emission density: market series in regime $k$ \\
$\xi_{t\mid\tau}(k)$ & $p(z_t{=}k\mid m_1,\dots,m_\tau)$: filtered ($\tau=t$), predicted ($\tau=t-1$), smoothed ($\tau=T$) \\
$\pi_k(t)=\mathbb P(z_{t+1}{=}k\mid\F_t)$ & the gate: weight of expert $k$ for the day-$(t{+}1)$ return \\
$\mu_k(x)=f_0(x)+c_k(x;\psi_k)$ & mean of expert $k$: frozen base plus correction \\
$s_k$ & noise standard deviation of expert $k$ \\
$\gamma$ & responsibility: posterior probability of a regime or expert after seeing the returns \\
\bottomrule
\end{tabularx}
\end{center}

\begin{warnbox}{Notation clashes to keep in mind}
\begin{itemize}
  \item Bishop writes the \emph{target} as $t$; here $t$ is always time.
  \item \path{Model_Derivation_BaseCase.md} uses $s_t$ for the regime, $\mathbf P$ for the transition matrix, $r_t$ for the market series, $\mu_k$ for the gate mean, and $r_k$ for the expert corrections. In these notes those are $z_t$, $A$, $m_t$, $\nu_k$ and $c_k$, which frees $r$ for returns and $\mu$ for expert means.
  \item $\sigma_k$ is always the \emph{gate's} regime volatility of the market series. $s_k$ is always the \emph{expert's} noise. They measure different things (Section 3.3).
  \item Chamroukhi \& Nguyen use $\pi$ both for the HMM initial distribution and for the logistic gate. They use $\alpha_k$ for constant mixing proportions, $\tau$ for cluster posteriors, $\gamma$ for regime posteriors, and $\xi$ for pairwise posteriors.
\end{itemize}
\end{warnbox}
\newpage
